% page 69 du fac-simile (image p-068 du scan)
% bloc 162.6 x 229.0 mm ; marge gauche 39.4 mm
\pgc{17.780mm}{0.000mm}{
\l{\cel{55}6l}
\l{}
\l{}
\l{\sou{beatifikar}. (\sou{trans.})\cel{2}(\sou{aludante la papo})\cel{2}Pozar inter la beati}
\l{\cel{3}(homo qua vivis sante). - DEFIRS.}
\l{}
\l{\sou{bechika}. Efikiva kontre tusado. - EFIS.}
\l{}
\l{\sou{bedo}. Bendo de tero quan onu garnisas ye flori, arbusti, en}
\l{\cel{3}gardeno. - DE.}
\l{}
\l{\sou{bedelo}. Adulto qua preiras o sequas la procesioni (generale}
\l{\cel{3}religiala). - DEFIRS.}
\l{}
\l{\sou{begino}.\cel{2}Sorto de regulierini qui vivis kongregacione segun}
\l{\cel{3}reguli monakeyala, sen vovir. - DEFIRS.}
\l{}
\l{\sou{begonio}. (bot.) Tipo di la familio \textquotedbl{}begoniacei\textquotedbl{}, quan onu}
\l{\cel{3}kultivas kom plezur-planto. - L.\cel{1}\sou{begonia}. - DEFIRS.}
\l{}
\l{\sou{beko.}\cel{1}I.Boko di la uceli, qua konsistas ye du lami korno-harda}
\l{\cel{3}(mandibuli) ed uzebla kom armo por atakar o por su defensar.}
\l{\cel{3}- II. La extremajo di ula objekti qui finas pintatre. - EFIS.}
\l{}
\l{\sou{bekafigo}. (\sou{zool.}) Ula mikra ucelo. L.\cel{1}\sou{sulvia hortensis}. - DEFIS.}
\l{}
\l{\sou{bekaso}. (\sou{zool}.) Ucelo migrera, ek la klaso \textquotedbl{}steltieri\textquotedbl{}, de la}
\l{\cel{3}familio \textquotedbl{}longa-beka\textquotedbl{}. - L.\cel{1}\sou{scolopax rusticola}. - DEFIRS.}
\l{}
\l{\sou{bela}.\cel{2}Qua vekigas la admiro-sentimento, per igar senteblano}
\l{\cel{3}per manifestar ta o ca perfektajo korpala od etikala. - EFI.}
\l{}
\l{\sou{beladono}. (\sou{bot.)}\cel{2}Planto herbacea, perena, ek la familio \textquotedbl{}so-}
\l{\cel{3}lanei\textquotedbl{}, venenoza en omna lua parti, e di qua la toxikeso}
\l{\cel{3}uzesas medicinale. - L.\cel{1}\sou{atropa belladonna}. - DEFIRS.}
\l{}
\l{\sou{belemnito}. Konko fosila, longa, cilindra, finanta per pinto.}
\l{\cel{3}- DEFIRS.}
\l{}
\l{\sou{belvedero}. Paviliono konstruktita en loko de ube la peizajo}
\l{\cel{3}esas agreabla e vasta, od an la somito di edifico de ube onu}
\l{\cel{3}vidas vaste. - DEFIR.}
\l{}
\l{\sou{bemolo}. (\sou{muziko}). Signo pozita avan muzik-noto por savigar ke}
\l{\cel{3}onu devas abasar ye un mi-tono ta noto ed olti di la sama}
\l{\cel{3}grado en la sama mezuro. - Signo an la klefo por savigar ke}
\l{\cel{3}onu devas abasar ye un mi-tono omna noti skribita sur la}
\l{\cel{3}sama lineo, en la muzikajo. - FIRS.}
\l{}
\l{\sou{bendo}.\cel{2}Areo limitizita per du linei paralela. - DEFIRS.}
\l{}
\l{\sou{benedikar}. (\sou{trans.}) Dicar vorti qui +wishas ad ulu lua feliceso.}
\l{\cel{3}- DEFIS.}
}
